\part{Monitoring, fault detection and self-healing}

\FigHeal

\section{Fault detection}

\begin{qa}{470}{What faults is CARS-SF designed to detect?}
Backhaul cut (loss of carrier); backhaul degradation (rising delay and loss); core switch failure (core unreachable
after failover); bottleneck congestion (drops); saturation (utilisation above 95\%); latency spikes (RTT above
200\,ms); total connectivity loss.
\end{qa}

\begin{qa}{471}{Which faults are explicitly represented in the test scenarios?}
\code{fiber_cut} (cut at $\frac13$, repair at $\frac23$), \code{backhaul_degradation} (+15\,ms and +0.3\% loss per
tick from $0.2$ of the run, cut at $0.7$), \code{core_switch_failure} (all core ports down at $\frac13$, repair at
$\frac23$). Congestion arises naturally from load.
\end{qa}

\begin{qa}{472}{How do you distinguish link failure, congestion, excessive latency, packet loss and telemetry failure?}
\begin{itemize}
\item Link failure: carrier transition on \code{s_edge-bh}.
\item Congestion: local drops, backlog or saturation.
\item Excessive latency: RTT above threshold; as a trend without local congestion, attributed to the backhaul.
\item Loss: degradation loss shows as lost probes and loss in captures.
\item Telemetry failure: \textbf{not distinguished}. An unavailable reader looks like zero (Q279).
\end{itemize}
\end{qa}

\begin{qa}{473}{Fault detection vs diagnosis vs localisation vs recovery?}
Detection: something is wrong. Diagnosis: what kind (cut vs degradation vs congestion). Localisation: where (backhaul
link vs core switch). Recovery: restoring service. CARS-SF localises by elimination: if failover to a path that
bypasses the link does not restore the core, the fault is the shared endpoint (\code{s_core}).
\end{qa}

\begin{qa}{474}{Is the detector rule-based, statistical, ML-based or hybrid?}
Hybrid rule and statistical: thresholds and link events (reactive), plus a least-squares RTT trend over a 6-sample
window with a horizon forecast (predictive). No ML.
\end{qa}

\begin{qa}{475}{Why is that detector appropriate here?}
\why{Faults must be acted on in under 1\,s, explainably, with no failure labels for training.}
\choice{Physical signals (carrier) for cuts; trends plus a congestion discriminator for degradation.}
\alt{An ML anomaly detector (isolation forest, LSTM) over telemetry.}
\trade{Deterministic, explainable and data-free; but thresholds are hand-set and novel fault types are missed.}
\ev{Detected every injected cut; no false failover in any run without an injected fault (Q487).}{\Lthree}
\end{qa}

\begin{qa}{476}{What are the inputs to the detector?}
The per-tick snapshot (drops delta, utilisation, per-service RTT, backhaul and backup flags, queue backlog) and
timestamped link-watcher events, delivered between ticks through the fast path.
\end{qa}

\begin{qa}{477}{What constitutes an anomaly?}
\begin{itemize}
\item More than 50 bottleneck drops per tick, or utilisation above 95\%.
\item RTT above 200\,ms.
\item All probes lost (connectivity loss).
\item Core probes lost with edge probes alive (core unreachable).
\item Carrier loss on the backhaul (cut).
\item Predicted degradation: every core-bound RTT slope at least 2\,ms per tick, the forecast 5 ticks ahead at least
the strictest core threshold, edge flat, and no local congestion in the window.
\end{itemize}
\end{qa}

\begin{qa}{478}{How are thresholds calibrated?}
By design and checked against scenario runs: the strictest core-bound threshold is twice education's 150\,ms SLA;
the slope floor of 2\,ms per tick lies above probe noise ($\approx$0.5\,ms).
\honest{No ROC calibration. The congestion discriminator was added after a real false failover on OVS.}
\end{qa}

\begin{qa}{479}{What are the trade-offs between false positives and false negatives?}
A false failover moves core traffic to a 50\,Mb/s, 40\,ms path, costing capacity and latency, and flaps. A missed cut
costs an outage. Cuts use a physical signal (near-zero false positives); prediction is conservative (several
conditions must hold), trading earlier warnings for fewer false alarms.
\end{qa}

\begin{qa}{480}{What if a genuine fault produces no immediate SLA violation?}
Carrier-based detection does not wait for an SLA violation: the cut is acted on at once. Degradation is caught as a
trend \emph{before} the SLA breaks (which is the point of prediction).
\end{qa}

\begin{qa}{481}{What if normal workload variation triggers an anomaly?}
Latency spikes and congestion in normal runs do trigger healthcare or affected-slice boosts. These are cheap,
bounded ($\times$2) and expire. Normal variation does \emph{not} trigger failover: \NFalseFailoverRuns{} false failovers
in \NNoFaultRuns{} runs without a fault.
\end{qa}

\begin{qa}{482}{How does the detector distinguish a transient spike from persistent degradation?}
Prediction needs a 6-sample trend with a consistent slope and no congestion anywhere in the window; escalation needs 2
consecutive ticks; boosts react to single ticks but expire.
\end{qa}

\begin{qa}{483}{What is the role of temporal windows and consecutive violations?}
They suppress noise: the trend window for prediction, 2 consecutive ticks for core-failure diagnosis, a 5-tick
hold-down before fail-back, and a 6-tick boost lifetime.
\end{qa}

\begin{qa}{484}{How do you detect a fiber cut if the interface stays present but traffic stops flowing?}
The interface stays present after a cut; what matters is carrier. On veth (as on fiber with loss of light), a far-end
failure sets our side to \code{lowerlayerdown}, which the watcher sees within 50\,ms.
\honest{A \emph{silent} failure (carrier up, forwarding dead) shows only as core-unreachable probes, and in the
\textsf{NORMAL} state that event does not trigger failover. Gap: add ``core unreachable for 2 ticks while NORMAL''
$\rightarrow$ fail over.}
\end{qa}

\begin{qa}{485}{How do you distinguish a dead path from a merely congested one?}
Dead: probes are lost, carrier is down, deliveries stop. Congested: probes return with high RTT, and local backlog and
drops rise.
\end{qa}

\begin{qa}{486}{How is diagnosis validated against known injected faults?}
Each resilience run records injector events (type, tick, timestamp) and healer actions (trigger, timestamp). The
report compares them: switched before the cut in \NDegBeforeCutShare{} of degradation runs; core-switch failure
diagnosed in \NCoreDiagnosedShare{}; an integration test asserts each.
\end{qa}

\begin{qa}{487}{How do you measure precision, recall, false-alarm rate and detection delay?}
\begin{itemize}
\item Recall: injected failures with a healing response, \NInjectedDetected/\NInjectedRuns.
\item False alarms: failovers in fault-free runs, \NFalseFailoverRuns/\NNoFaultRuns.
\item Detection delay: MTTD \NMttdCutFull\,s for a cut (event-driven); \NMttdDegFull\,s from the start of degradation
for prediction (by design it waits for the trend).
\end{itemize}
\honest{Small $n$ per fault type (3 seeds $\times$ 3 variants), so no confidence intervals on recall.}
\end{qa}

\begin{qa}{488}{How does detector performance change with several simultaneous faults?}
\honest{Untested. Combinations (degradation + congestion, cut + context change) are not in the matrix. Congestion
during degradation would, by design, suppress prediction (conservative); cuts would still be caught.}
\end{qa}

\begin{qa}{489}{What happens with a fault outside the detector's known categories?}
It is caught only if it shows up as one of the signals (RTT, drops, carrier, reachability). Otherwise it is
invisible. The SLA monitor still records the violation.
\end{qa}

\begin{qa}{490}{How does the system avoid acting on unreliable or stale observations?}
Link events carry timestamps; self-inflicted transitions (within 1\,s of the healer's own action) are ignored;
prediction requires sustained trends.
\honest{There is no general staleness check (Q280).}
\end{qa}

\section{Recovery planning and action selection}

\begin{qa}{491}{What recovery actions can CARS-SF execute?}
\begin{itemize}
\item Quarantine the primary backhaul, activate the backup, flush the FDBs.
\item Fail back (backup down, then primary up, then flush).
\item Allocation constraints: core slices capped to 95\% of backup capacity with HTB ceilings; degraded mode holds
core slices at their floors and lends the rest to edge slices.
\item Priority boosts (affected slice $\times$2; general $\times$0.5 on saturation).
\item Diagnosis escalation.
\end{itemize}
\end{qa}

\begin{qa}{492}{Which actions change bandwidth allocation, and which change the forwarding path?}
Path: quarantine, backup activation, fail-back. Allocation: constraints, ceilings, boosts.
\end{qa}

\begin{qa}{493}{How is the candidate action set constructed?}
A fixed mapping from failure type to action (state machine). There is no search over candidates.
\end{qa}

\begin{qa}{494}{How does the planner decide which action to attempt first?}
Path protection first (it restores connectivity), then allocation constraints (the controller re-plans at once because
the healing state changed).
\end{qa}

\begin{qa}{495}{Why does context influence recovery ordering?}
It does not influence the \emph{ordering}; it influences the \emph{post-recovery allocation}, which uses the current
context's weights under the healing constraints.
\honest{Context-dependent recovery (e.g.\ preferring a costlier path only in an emergency) is not implemented.}
\end{qa}

\begin{qa}{496}{How does the system account for resource capacity when choosing an action?}
The constraints encode the backup's capacity (core slices $\le$0.95$\times$50\,Mb/s) and degraded mode. Without the
ceilings, the 50\,Mb/s backup bufferbloated (RTT oscillating 116--336\,ms); with them it stays at its 96\,ms floor.
\end{qa}

\begin{qa}{497}{What if the preferred recovery action is infeasible?}
If the backup does not restore the core, the engine escalates to \textsf{DEGRADED} (graceful degradation). There is
no other path to try.
\end{qa}

\begin{qa}{498}{How does the planner compare predicted outcomes of candidate actions?}
\honest{It does not. There is no outcome model in the planner.}
The calibrated sequential simulator could score candidate allocations before applying them.
\end{qa}

\begin{qa}{499}{What evidence supports the prediction model?}
For the \emph{fault predictor}: switched before the cut in \NDegBeforeCutShare{} of runs; outage \NOutDegFullMean{} vs
\NOutDegReactiveMean\,s reactive-only. There is no outcome-prediction model for actions.
\end{qa}

\begin{qa}{500}{What happens when the predicted outcome differs from the actual outcome?}
For prediction, the guards handle it: if the primary later proves healthy (OAM), fail-back follows after hold-down.
A wrong failover costs only capacity until then.
\end{qa}

\begin{qa}{501}{How do you prevent an action from degrading another service's QoS?}
Floors always hold; constraints only move capacity from slices that cannot use it (core-bound slices during an outage)
to edge slices; boosts are bounded and expire.
\end{qa}

\begin{qa}{502}{How do you avoid oscillating between two allocations or paths?}
A 5-tick hold-down; OAM-gated fail-back; ignoring self-inflicted link events; the change threshold; boost cooldowns.
\ev{Before these guards, a degradation run alternated failover and fail-back 3 times in 3 ticks. After them, exactly
one proactive switchover. An OVS run with AQM-less queues flapped before the congestion discriminator and did not
after.}{\Lthree}
\end{qa}

\begin{qa}{503}{Is there a maximum number of recovery attempts?}
Path switching happens once per failure episode (state-guarded).
\honest{Boosts can recur every 6 ticks indefinitely; there is no attempt counter.}
\end{qa}

\begin{qa}{504}{How are retry delays and back-off selected?}
Fixed: a 6-tick boost lifetime and cooldown, a 5-tick hold-down, 2 ticks before escalation. Chosen to exceed the
5\,s control interval and TCP's reaction time.
\honest{There is no exponential back-off.}
\end{qa}

\begin{qa}{505}{How does the planner tell a failed action from one that needs more time?}
After backup activation it waits 2 ticks (and requires 2 consecutive core-unreachable ticks) before concluding the
core itself is down.
\end{qa}

\begin{qa}{506}{Reactive recovery vs proactive reconfiguration?}
Reactive: after the cut (carrier). Proactive: before the cut, when trends predict a breach. Measured outage: proactive
\NOutDegFullMean\,s vs reactive \NOutDegReactiveMean\,s on degradation scenarios.
\end{qa}

\begin{qa}{507}{When should the system escalate to an operator?}
On \textsf{DEGRADED} (core-switch failure, which needs physical repair), on repeated boosts with no SLA recovery, on
configuration failures, and on any state that persists beyond a time limit.
\honest{Escalation is only logged; there is no paging integration.}
\end{qa}

\begin{qa}{508}{How do you ensure the controller does not endlessly repeat an ineffective action?}
Path actions are state-guarded (one per episode).
\honest{Boost repetition is not limited. Add an attempt budget and escalation after N ineffective boosts.}
\end{qa}

\FigResilience

\section{The closed-loop control system}

\begin{qa}{509}{What are the precise stages of the CARS-SF self-healing loop?}
Monitor (telemetry, SLA monitor, link watcher) $\rightarrow$ Analyse (reactive detector, predictor) $\rightarrow$
Plan (healing state machine, controller re-plan under constraints) $\rightarrow$ Execute (topology primitives, QoS
backend) $\rightarrow$ Knowledge (registries and history). This is a MAPE-K loop.
\end{qa}

\begin{qa}{510}{What is the role of monitoring in the loop?}
To provide the observations (rates, drops, RTT, carrier) and the per-slice SLA record against which success is
judged.
\end{qa}

\begin{qa}{511}{What is the role of diagnosis?}
To pick the right action: a cut means fail over; a core failure means degrade gracefully (failover is pointless);
congestion means re-prioritise and do \emph{not} fail over.
\end{qa}

\begin{qa}{512}{What is the role of the recovery planner?}
To map the diagnosis to a safe action sequence (loop-free ordering) and to constraints for the allocator.
\end{qa}

\begin{qa}{513}{What is the role of the enforcement backend?}
To execute path changes (\code{ip link}, FDB flush) and rate changes (tc, OVS or UPF).
\end{qa}

\begin{qa}{514}{What is the role of post-action verification?}
To confirm that the action achieved the goal before declaring success, or to trigger the next step. CARS-SF verifies
\emph{reachability} after failover (core-unreachable escalation) and \emph{link health} before fail-back (far end up
plus OAM).
\honest{There is no generic KPI check after allocation changes and no rollback.}
\end{qa}

\begin{qa}{515}{Why is verification necessary after a recovery action?}
Actions can fail silently (a dead backup, a tc error, stale FDB entries) or fix the wrong problem (failover when the
core switch is the fault). The core-switch case is exactly where verification changed the decision.
\end{qa}

\begin{qa}{516}{What evidence is required to declare recovery successful?}
Online: core probes answered again on the new path. Post-run (ground truth): the measured outage (longest delivery
gap) and packets lost. Fiber cut: \NOutCutFullMean\,s mean outage.
\end{qa}

\begin{qa}{517}{How long should the system wait for the network to settle?}
At least the FDB relearning time plus TCP recovery, about 1\,s on OVS. We use 2 ticks before escalation and 5 ticks of
hold-down before fail-back.
\end{qa}

\begin{qa}{518}{How many consecutive good measurements before declaring recovery?}
Implicitly 2 (escalation logic) for the core path.
\honest{There is no explicit ``recovered'' declaration with N good SLA samples. Add one with N=3.}
\end{qa}

\begin{qa}{519}{What if the SLA improves temporarily and then degrades again?}
On backup, renewed degradation would show in probes. Escalation and boosts react. A repeated primary failure after
fail-back starts a new episode.
\end{qa}

\begin{qa}{520}{What if the action improves latency but drops throughput below the SLA?}
That is what happens on backup: capacity halves, so core slices are capped and their satisfaction drops while
connectivity is restored. The SLA monitor records it. It is a deliberate trade-off (connectivity over throughput),
reported per slice.
\end{qa}

\begin{qa}{521}{How do you evaluate recovery when several KPIs must hold at once?}
Per-slice compliance requires latency, loss and satisfaction together (\code{sla_ok}). Run-level: outage, packets lost
and weighted utility (\NOutCutFullMean\,s, \NLostCutFull{} packets, utility in \code{resilience_summary.csv}).
\end{qa}

\begin{qa}{522}{What if the verifier reports that recovery failed?}
For a failover that did not restore the core: escalate to \textsf{DEGRADED}. For fail-back with OAM failing: stay on
backup.
\end{qa}

\begin{qa}{523}{Does the controller roll back, escalate, or try another action?}
It escalates (degraded mode).
\honest{It does not roll back to the primary while the primary is down, and there is no alternative action list
beyond that.}
\end{qa}

\begin{qa}{524}{How do you ensure rollback does not worsen the original fault?}
Fail-back is the rollback, and it is gated by far-end-up plus OAM health plus hold-down, so we never return to a
known-bad link.
\end{qa}

\begin{qa}{525}{How is every decision, action and verification result recorded?}
\code{healing_actions} (action, tick, timestamp, trigger, trigger timestamp, state); \code{injector_events};
\code{failure_types}; \code{controller_history}; \code{feedback.csv} (per-tick healing state); \code{resilience}
summary, all in \code{summary.json}.
\end{qa}

\begin{qa}{526}{Open-loop recovery script vs verified closed-loop controller?}
An open-loop script executes a fixed action. A closed-loop controller observes the outcome and adapts (escalation,
OAM-gated fail-back, hold-down). CARS-SF is closed-loop for path recovery, but only partially verified for allocation
changes.
\end{qa}

\begin{qa}{527}{What experiment proves that verification reduces unnecessary reconfigurations?}
Anecdotal but measured: without the hold-down, OAM gate and self-event suppression, a degradation run produced 3
failover/fail-back cycles; with them, 1 switchover.
\honest{There is no controlled ablation across seeds. Add a ``no guards'' variant to the resilience study.}
\end{qa}

\begin{qa}{528}{How would you demonstrate the entire loop in a live Mininet run?}
\code{sudo .venv/bin/python run_experiment.py --scenario backhaul_degradation --duration 45 --apply-qos}. Watch the
log: degradation start; predicted degradation; proactive backup activation; ceilings applied; the cut (no outage);
fail-back after repair. Then open \code{summary.json}: \code{resilience}, \code{healing_actions}.
\end{qa}

\section{Fault tolerance of the controller itself}

\begin{qa}{529}{What happens if the monitoring component crashes?}
It runs in the orchestrator process. A crash ends the run; the kernel keeps the last QoS state.
\honest{There is no supervisor, so a crash means loss of healing.}
\end{qa}

\begin{qa}{530}{What happens if the controller loses contact with the switch?}
Not applicable in emulation (local commands). In a deployment it would be a partition: enforcement persists, healing
on the far side stops (Q157).
\end{qa}

\begin{qa}{531}{What happens if the controller restarts mid-recovery?}
\honest{State is in memory only. After a restart the engine starts in \textsf{NORMAL} while the primary is quarantined
and the backup active.}
Fix: persist the healing state, and on start reconcile from link admin states.
\end{qa}

\begin{qa}{532}{How is the previous network configuration recovered after a crash?}
It is not reconstructed. The kernel retains qdiscs and link states, and Mininet teardown in \code{finally} cleans up
the emulation. A real deployment would read back state on start.
\end{qa}

\begin{qa}{533}{How do you prevent two controller instances from issuing conflicting actions?}
\honest{Nothing does. A single instance is assumed.} Use a lease or leader election (etcd or ZooKeeper) and
idempotent intents.
\end{qa}

\begin{qa}{534}{How do you handle duplicate commands and retries?}
Most commands are idempotent (\code{ip link set up/down}, \code{tc class change} to absolute values), and the injector
skips duplicates. \code{tc filter add} is not idempotent, so flow rules are tracked and installed once.
\end{qa}

\begin{qa}{535}{Are configuration changes idempotent?}
Rates and ceilings, link states and FDB flushes are idempotent. Filter additions are guarded by the flow-rule table.
\end{qa}

\begin{qa}{536}{How do you ensure the system eventually converges to a stable state?}
Finite states with guarded transitions (hold-down, OAM, one switch per episode); expiring boosts; the change
threshold. In tests and runs it always converged.
\honest{There is no formal convergence proof.}
\end{qa}

\begin{qa}{537}{What if the primary and backup paths fail simultaneously?}
That is equivalent to the core-switch case: escalate to \textsf{DEGRADED}, protect edge services, wait for repair.
\end{qa}

\begin{qa}{538}{What if no feasible recovery action exists?}
Degrade gracefully (keep surviving services prioritised), keep monitoring, escalate to the operator.
\end{qa}

\begin{qa}{539}{How does the controller preserve critical traffic during partial failure?}
Floors, context weights, and degraded-mode reallocation: core slices are held at their floors and edge slices,
including healthcare (served by the edge), get the capacity.
\end{qa}

\begin{qa}{540}{What is the minimum safe behaviour when the controller cannot decide reliably?}
Keep the last good configuration, never form loops, never fail back without OAM, keep floors, alert the operator.
The kernel's persistence of HTB state provides much of this for free.
\end{qa}
